\chapter{پیاده‌سازی سامانه}

\section{ساختار مخزن}
مخزن پروژه یک پروژه چندماژوله است. فایل اصلی \lr{\texttt{pom.xml}} ماژول‌های \lr{\texttt{event-receiver}}، \lr{\texttt{kernel-event-receiver}}، \lr{\texttt{event-aggregator}} و \lr{\texttt{profiler-evaluation}} را تعریف می‌کند و نسخه جاوا را 17 قرار می‌دهد. بخش \lr{\texttt{event-analytics}} با پایتون پیاده‌سازی شده و وابستگی‌های آن در فایل \lr{\texttt{requirements.txt}} آمده است. بخش \lr{\texttt{kernel-event-sender}} نیز شامل کد \lr{eBPF}، بارگذار \lr{C} و \lr{Makefile} است.

\section{دریافت رخدادهای \lr{JFR}}
ورودی اصلی رخدادهای سطح برنامه در کلاس \lr{\texttt{aut.ce.App}} از ماژول \lr{\texttt{event-receiver}} آغاز می‌شود. این کلاس نام ماشین مجازی هدف را از متغیر محیطی \lr{\texttt{TARGET\_VM\_NAME}} می‌خواند و در صورت نبود مقدار، نام \lr{\texttt{JfrEventGenerator}} را در نظر می‌گیرد. سپس \lr{\texttt{JfrRemoteListener}} و \lr{\texttt{JfrRemoteEventStreamer}} اجرا می‌شوند و نخ \lr{\texttt{KafkaWriter}} پیام‌ها را به \lr{Kafka} ارسال می‌کند.

برای استخراج فیلدهای تخصصی، الگوی غنی‌ساز رخداد استفاده شده است. هر غنی‌ساز بخشی از بار رخداد را می‌سازد و خروجی نهایی به صورت پیام ساخت‌یافته ارسال می‌شود. این روش باعث می‌شود اضافه کردن یک نوع رخداد جدید با افزودن یک غنی‌ساز جدید انجام شود.

\section{دریافت رخدادهای هسته}
بخش \lr{\texttt{kernel-event-sender}} شامل برنامه \lr{eBPF} و بارگذار کاربرفضا است. برنامه \lr{eBPF} رخدادهای مرتبط را از هسته دریافت می‌کند و بارگذار آن‌ها را به مقصد شبکه ارسال می‌کند. ماژول \lr{\texttt{kernel-event-receiver}} این رخدادها را دریافت، بافر و در \lr{Kafka} منتشر می‌کند. این مسیر برای سناریوهایی مفید است که نشانه اصلی آن‌ها در سطح سیستم‌عامل دیده می‌شود.

\section{تجمیع رخدادها}
کلاس \lr{\texttt{EventAggregatorApp}} قلب ماژول تجمیع است. این کلاس دو توپولوژی مشابه می‌سازد: یکی برای رخدادهای \lr{JFR} و دیگری برای رخدادهای هسته. رخدادها ابتدا فیلتر می‌شوند تا پیام‌های نامعتبر حذف شوند. سپس بر اساس کلید تجمیع گروه‌بندی شده و در پنجره‌های زمانی پردازش می‌شوند. خروجی هر پنجره شامل زمان شروع و پایان، مدت پنجره، کلید و معیارهای محاسبه‌شده است.

راهبردهای تجمیع در بسته \lr{\texttt{strategy}} قرار دارند. این راهبردها تفاوت رخدادهای پردازنده، حافظه، فایل، شبکه و جمع‌آوری زباله را از منطق عمومی \lr{Kafka Streams} جدا می‌کنند. در نتیجه، تغییر روش محاسبه یک معیار نیازمند تغییر توپولوژی اصلی نیست.

\section{تحلیل و گزارش‌گیری}
در ماژول \lr{\texttt{event-analytics}}، فایل \lr{\texttt{main.py}} پیکربندی را می‌خواند، مصرف‌کننده \lr{Kafka} را می‌سازد، آشکارسازها را مقداردهی می‌کند و پیام‌های تجمیع‌شده را پردازش می‌کند. آشکارسازهای موجود شامل پردازنده، حافظه، ورودی/خروجی فایل، شبکه، قفل و هسته هستند. خروجی‌ها توسط \lr{\texttt{CsvReportWriter}} در قالب گزارش رخدادهای تجمیع‌شده و گزارش ناهنجاری نوشته می‌شوند.

\section{پایش خود سامانه}
پروژه برای مشاهده وضعیت اجزای خود نیز از \lr{Prometheus} استفاده می‌کند. در ماژول‌های جاوا، کلاس‌های \lr{\texttt{MetricsHandler}} معیارهایی مانند مقدارهای تجمیع‌شده را به‌روزرسانی می‌کنند. پوشه \lr{\texttt{monitoring}} شامل پیکربندی لازم برای اجرای \lr{Prometheus} است \cite{prometheusdocs}.

\section{جمع‌بندی}
پیاده‌سازی پروژه بر اصل جداسازی مسئولیت‌ها تکیه دارد. دریافت‌کننده‌ها مسئول تبدیل رخداد خام به پیام، تجمیع‌کننده مسئول تبدیل پیام‌ها به معیارهای پنجره‌ای، و تحلیل‌گر مسئول تصمیم‌گیری درباره ناهنجاری‌ها است.
